% ======================================================================
\section{The chemical solve}
\label{sec:chemistry}
% ======================================================================

\subsection{Network and rate law}
\label{sec:network}

The network is 3D-PDR's REDUCED network (files
\src{dataPDR/\allowbreak{}species\_reduced.d} and \src{dataPDR/\allowbreak{}rates\_reduced.d}), in the
format of the UMIST Database for Astrochemistry. \citet{Bisbas2012}
introduced its 33 species with 320 reactions drawn from UMIST 2006
\citep{Woodall2007}; \citet{Bisbas2015} describe the 33-species network then in
use as a 330-reaction subset of UMIST 2012 \citep{McElroy2013}. The vendored
file, from upstream commit \texttt{fabcdc6}, has 331 reactions and records
neither its UMIST release nor the references of its coefficients, so the
release cannot be determined from the files. It has 33 species,
\begin{center}
\small H, H$_2$, H$^+$, H$_2^+$, H$_3^+$, He, He$^+$, C, C$^+$, CH, CH$^+$,
CH$_2$, CH$_2^+$, CH$_3$, CH$_3^+$, CH$_4$, CH$_4^+$, CH$_5^+$, O, O$^+$, OH,
OH$^+$, H$_2$O, H$_2$O$^+$, H$_3$O$^+$, O$_2$, O$_2^+$, CO, CO$^+$, HCO$^+$, Mg,
Mg$^+$, e$^-$,
\end{center}
and 332 reactions: 3D-PDR's 331 plus reaction 332,
$e^-+\mathrm{H_2^+}\to e^-+\mathrm H+\mathrm H^+$ ($2\times10^{-7}e^{-30720/T}$\,cm$^3$\,s$^{-1}$),
appended in this project because the network had no
electron-impact dissociation of \Htwo$^+$, Cloudy's main \Htwo$^+$ sink in ionised
gas (\src{pdr/\allowbreak{}UPSTREAM.md}, \dlog{D-395}); the Arrhenius form is the one of
Cloudy's \src{mole\_co\_umisthack.dat}, which attributes it to
\citet{Janev1987}. Of the 332, 28 are photoreactions, 8 cosmic-ray and 15
cosmic-ray-induced photoreactions (Section~\ref{sec:photorates}), and three
carry the grain marker \texttt{\#}: grain formation of \Htwo{} (328) and the
grain-assisted recombinations of He$^+$ and C$^+$ (330, 331). These two have no
term in the generated right-hand side, so their coefficients are computed but
unused and 330 reactions act in the solve. The other 278 are two-body reactions:
ion--molecule, neutral--neutral and charge-transfer reactions, dissociative and
radiative recombination, radiative association and collisional dissociation.
The ODE right-hand side is generated C code selected at compile time
(\src{pdr/\allowbreak{}odes\_reduced.c}, flag \kw{-DREDUCED}); the keyword \kw{pdrNetwork}
only validates that the requested network is the compiled one.

Two-body coefficients follow the modified Arrhenius form
\begin{equation}
  k(T) = \alpha\left(\frac{T}{300\,\mathrm K}\right)^{\beta} e^{-\gamma/T}\quad[\mathrm{cm^3\,s^{-1}}],
  \label{eq:arrhenius}
\end{equation}
with $\gamma$ in K (\src{pdr/\allowbreak{}calc\_reac\_rates.F90:\allowbreak{}172--252}). Where a reaction
has several entries for different temperature ranges, the first whose upper
limit is above $T$ (else the last) is used and the others are set to zero; a
fit with $\gamma<0$ is evaluated at the lower limit of its range below it, and one
with $\gamma<-200$\,K is set to zero there. The chemistry is evaluated at
MOCASSIN's $T_e$ clamped to [\kw{pdrTmin}, \kw{pdrTmax}] $=[1,\,2\times10^4]$\,K by
default (\src{pdr\_mod.f90:\allowbreak{}3346}), the range over which the fits are meant to be
used; MOCASSIN's own floor \kw{teFloor}, 2.725\,K by default, keeps $T_e$ above
$T_{\mathrm{CMB}}$. The vendored bound $\max(\kw{pdrTmin},T_{\mathrm{CMB}})$
(\src{:\allowbreak{}364}) is read only by 3D-PDR's own temperature iteration
(\src{pdr/\allowbreak{}changetemperature.F90}), which the bridge does not call. In terms of fractional abundances $x_i=n_i/n_{\mathrm H}$ the system is
\begin{equation}
  \frac{\dd x_i}{\dd t} = \sum_{j,k} k_{jk}\,x_jx_k\,n_{\mathrm H}
      + \sum_j k_j\,x_j - x_i\Bigl(\sum_j k_{ij}\,x_j\,n_{\mathrm H} + \sum k_i\Bigr),
  \label{eq:odes}
\end{equation}
with two-body coefficients $k_{jk}$, $k_{ij}$ [cm$^3$\,s$^{-1}$] and first-order rates
$k_j$, $k_i$ [s$^{-1}$] (photoreactions and cosmic rays, Section~\ref{sec:photorates}).
Grain formation enters as $R_{\mathrm f}\,x(\mathrm H)\,n_{\mathrm H}$ in $\dd x(\mathrm{H_2})/\dd t$
and twice that as a loss of H. Inside the right-hand side the electron
abundance is not a free variable: it is the sum of the network's cations plus
\rtn{xe\_extra}, the electrons of the ions the network has no species for
(higher stages, other elements), supplied by the bridge (\dlog{D-402};
\src{pdr/\allowbreak{}odes\_reduced.c:\allowbreak{}58}).

\subsection{Initial conditions and seeding}
\label{sec:seed}

\paragraph{First visit of a cell.} \rtn{seedAbunFromMocassin}
(\src{pdr\_mod.f90:\allowbreak{}3123--3157}, \dlog{D-208}): every molecule starts at zero;
each linked element's neutral and first ion take MOCASSIN's stage-1 and stage-2
fractions times the element abundance; $x_e=n_e/n_{\mathrm H}$. Nothing is read
from the network file's dark-cloud defaults.

\paragraph{First PDR iteration.} \rtn{seedMolecules}
(\src{pdr\_mod.f90:\allowbreak{}2855--3023}, \dlog{D-209}) places a first \Htwo{} and CO
content so that the first shielding columns are not zero. \Htwo{} comes from the
balance of grain formation and unshielded photodissociation, with the neutral
hydrogen of MOCASSIN, $x_{\mathrm n}n_{\mathrm H}$, shared between H and \Htwo:
\begin{gather}
  R_{\mathrm f}\,n_{\mathrm H}\,n(\mathrm H) = D\,n(\mathrm{H_2}),\qquad
  n(\mathrm H)+2n(\mathrm{H_2}) = x_{\mathrm n}n_{\mathrm H}
  \quad\Longrightarrow\quad
  x(\mathrm{H_2}) = \frac{R_{\mathrm f}\,x_{\mathrm n}\,n_{\mathrm H}}{D+2R_{\mathrm f}n_{\mathrm H}},\\
  D = \alpha_{\mathrm{H_2}}\,\chi\;\mathrm{SCATTER}(\AV,1000\,\text{\AA}),
\end{gather}
with $\alpha_{\mathrm{H_2}}$ the network's coefficient ($5.18\times10^{-11}$\,s$^{-1}$, read
from the rate file) and $\chi$ the local MC field of Equation~(\ref{eq:chi});
$x(\mathrm{H_2})$ is capped at $x_{\mathrm n}/2$. CO follows the ansatz
$x(\mathrm{CO}) = x(\mathrm C)\,[2x(\mathrm{H_2})/x_{\mathrm n}]$, capped by the oxygen,
with carbon and oxygen conserved. Steps 4--5 of Table~\ref{tab:steps} are then
repeated on the seeded abundances. The seed attenuates the already attenuated
local field $\chi$ once more with \rtn{SCATTER} and ignores self-shielding; it
only provides the starting point, which the iteration then forgets.

\subsection{Imposition of MOCASSIN's ionisation before the solve}
\label{sec:impose}

MOCASSIN owns the ionisation state; the chemistry owns what is locked in
molecules (\dlog{D-153}, \dlog{D-193}, \dlog{D-402};
\src{pdr\_mod.f90:\allowbreak{}3503--3584}). For each linked element $X$ with abundance
$A_X$ and molecular content $M_X=\sum_i \nu_{X,i}x_i$ (atoms of $X$ in each
species, from the species' formulae, \rtn{speciesAtoms}):
\begin{equation}
  x(X) = (A_X-M_X)\,\frac{f_1}{\sum_s f_s},\qquad
  x(X^+) = (A_X-M_X)\,\frac{f_2}{\sum_s f_s},
\end{equation}
where $f_s$ are MOCASSIN's stage fractions; what MOCASSIN holds in stages
above the second stays MOCASSIN's. The electron abundance is set to
$n_e/n_{\mathrm H}$ from MOCASSIN, and \rtn{xe\_extra} to $x_e$ minus the
linked first ions (\src{:\allowbreak{}3565--3581}). The imposition fixes the starting
point of the integration; inside it, the linked atoms are ionised and
recombined at MOCASSIN's own rates (Section~\ref{sec:linked}), so the solver's
neutral/ion ratio should stay close to MOCASSIN's, and the remaining
difference is carried by the network's other channels (charge transfer,
ion--molecule reactions).

\subsection{Integration to steady state}

\rtn{calculate\_abundances} (\src{pdr/\allowbreak{}calculate\_abundances.c}), called at
\src{pdr\_mod.f90:\allowbreak{}3749}, integrates Equation~(\ref{eq:odes}) with CVODE from
SUNDIALS 7.8.0 (\dlog{D-015}; \citealt{Hindmarsh2005,Gardner2022}). The line numbers
in this list are those of \src{pdr/\allowbreak{}calculate\_abundances.c} unless a file is
named:
\begin{itemize}
\item variable-order backward differentiation formulae (\verb|CV_BDF|, :196),
  the stiff integrator, with Newton iteration on a dense direct linear solver
  (\verb|SUNDenseMatrix|, \verb|SUNLinSol_Dense|, :229--238) and CVODE's
  difference-quotient Jacobian (the analytic Jacobian is not used);
\item relative and absolute tolerances $10^{-8}$ and $10^{-30}$ on the fractional
  abundances (set at \src{pdr\_mod.f90:\allowbreak{}369--370}, applied at :212); at most
  $10^7$ internal steps per output interval (:115);
\item output times from $10^{-4}$\,yr, multiplied by 10 until
  \kw{pdrEndTime} (default $10^7$\,yr) (:246, :364);
\item \kw{CHEMRETRY} (:248--384): a failed integration is restarted from the
  initial state up to four more times with a forced first step of
  $10^{-2}$\,s, divided by 100 at each retry;
\item \kw{CHEMSTEADY} (:387--437): the integration then continues in blocks that
  each double the elapsed time, at most seven ($128\times$\kw{pdrEndTime},
  $1.3\times10^9$\,yr by default), until no species above the absolute tolerance
  changes by more than 1\% over a block;
\item abundances below the absolute tolerance are set to zero and the
  electron entry is restored (:440--452).
\end{itemize}
The result is meant to be the steady state of the network at fixed
temperature, field, columns and imposed ionisation; there is no time
dependence carried between MC iterations. A species that relaxes
exponentially on a time scale $\tau$ from an offset of order unity passes the 1\%
test over a doubling once the elapsed time exceeds about $4.6\,\tau$, so species
that relax more slowly than about $10^8$\,yr cannot pass it by the last block;
if the seventh block does not pass, the abundances of that time are returned
without a flag. An
unchanged abundance vector is counted (\rtn{pdrNstuck},
\src{pdr\_mod.f90:\allowbreak{}3860}) but not treated as a failure, since it is also the
signature of a converged cell. The slowest chemical process in molecular gas
is usually \Htwo{} formation, whose time scale $1/(R_{\mathrm f}n_{\mathrm H})$ is
$\sim10^7$\,yr at $n_{\mathrm H}=100$\,cm$^{-3}$; where the age of the object is shorter
the steady state overstates the molecular content
(Section~\ref{sec:assumptions}).

\subsection{After the solve}

\begin{enumerate}
\item The solver's own neutral/ion split of each linked element is
  overwritten with MOCASSIN's stage-1:stage-2 ratio, keeping the solver's
  total $x(X)+x(X^+)$ (\src{pdr\_mod.f90:\allowbreak{}3770--3784}, \dlog{D-162}); the
  electron abundance is reset to MOCASSIN's (\src{:\allowbreak{}3801--3804}, \dlog{D-231}).
  The solver's own ratio is kept for diagnostics (\dlog{D-402}, \dlog{D-404}).
\item The abundances are stored (\src{:\allowbreak{}3922}) and $n(\mathrm{H_2})=x(\mathrm{H_2})n_{\mathrm H}$
  is written for MOCASSIN's opacity (\src{:\allowbreak{}4573}).
\item What the network does to each linked element's $X$ and $X^+$ beyond the
  channels \rtn{ionBalance} computes itself (photoionisation, cosmic rays,
  radiative recombination, charge transfer with H) is classified reaction by
  reaction at the solver's state, as net conversions $X\leftrightarrow X^+$ and
  as sources and sinks to and from molecules, and stored for MOCASSIN's next
  ionisation solve (\rtn{pdrIonRate}, \src{:\allowbreak{}3924--4054}, \dlog{D-402}). The
  classified totals reproduce the right-hand side's own formation and
  destruction totals to $4\times10^{-16}$ (\dlog{D-402}).
\end{enumerate}
